\chapter{\MakeUppercase{Fundamentação Teórica}}\label{sec:fundamentacao}

Esta seção apresenta o modelo do WSP, os métodos de solução e as técnicas usadas para estudar as instâncias
e a busca. A análise do espaço de instâncias (ISA), as medidas topológicas
e as redes de trajetórias de busca (STNs) fornecem a base para relacionar
propriedades dos problemas ao desempenho dos algoritmos.

\section{Problema de Supressão de Incêndios Florestais}\label{sec:wsp}

O WSP integra propagação do fogo e alocação temporal de recursos sobre um
grafo direcionado. A formulação com liberações de recursos ao longo do tempo
foi introduzida por \citeonline{Alvelos2018} e posteriormente resolvida por
métodos exatos e heurísticos
\cite{DelazeriRitt2024,HarrisForbesTaimre2023,MendesAlvelos2023}.

Seja $G=(V,A)$ um grafo direcionado cujos vértices representam células de uma paisagem
discretizada e cujos arcos indicam direções possíveis de propagação.
Cada arco $(u,v)$ possui tempo positivo $t_{uv}$, que pode incorporar
vento, inclinação e propriedades do combustível. Não se exige
$t_{uv}=t_{vu}$: a direção e a assimetria dos tempos participam do modelo.
Dado um vértice de ignição $s$, o fogo começa no instante zero e, sem
supressão, alcança $v$ no instante:
\begin{equation}
 a_v=d_G(s,v)=\min_{p\in\mathcal{P}_{sv}}\sum_{(u,w)\in p}t_{uw},
 \label{eq:tempo-chegada-livre}
\end{equation}
\noindent em que $\mathcal{P}_{sv}$ contém os caminhos direcionados de $s$
a $v$. Adota-se $a_v=+\infty$ quando $v$ é inalcançável. As intervenções de supressão, representadas pela aplicação de recursos
para retardar a propagação, podem modificar os caminhos que determinam esses tempos
(\citeauthor{DelazeriRitt2024}, \citeyear{DelazeriRitt2024,DelazeriRitt2026}).

Considere $k$ recursos, identificados por $R=\{1,\ldots,k\}$, com instantes
de liberação $r_j$. Uma alocação parcial é uma função injetiva
$\Lambda:R'\rightarrow V$, em que $R'\subseteq R$ reúne os recursos
efetivamente empregados. Cada recurso é associado a um vértice, e cada
vértice recebe no máximo um recurso. O conjunto protegido
$P^{\Lambda}=\operatorname{im}(\Lambda)$ é o conjunto dos vértices
que receberam essas alocações. No modelo básico,
cada recurso acrescenta o mesmo atraso $\Delta\geq0$ aos arcos de saída
do vértice que o recebe:
\begin{equation}
 t_{uv}^{\Lambda}=t_{uv}+\Delta\,\mathbb{I}[u\in P^{\Lambda}],
 \qquad (u,v)\in A.
 \label{eq:tempo-propagacao-alocacao}
\end{equation}
\noindent A função indicadora $\mathbb{I}[u\in P^{\Lambda}]$ vale um quando
$u$ pertence ao conjunto protegido e zero caso contrário. Assim, o atraso
é acrescentado apenas aos arcos que saem de vértices protegidos. Os novos
tempos de chegada são denotados por $a_v^{\Lambda}$. A viabilidade temporal exige
$a_{\Lambda(j)}^{\Lambda}\geq r_j$ para todo recurso empregado
(\citeauthor{DelazeriRitt2024}, \citeyear{DelazeriRitt2024,DelazeriRitt2026}). Portanto, proteger um vértice
significa retardar a propagação, sem necessariamente torná-lo imune.
O benefício pode ocorrer em vértices posteriores, e caminhos alternativos
podem passar a ser os mais rápidos.

Para um horizonte $H$, o objetivo é minimizar o número de vértices
atingidos estritamente antes dele:
\begin{equation}
 \min_{\Lambda\ \text{viável}} |B_H^{\Lambda}|,
 \qquad B_H^{\Lambda}=\{v\in V:a_v^{\Lambda}<H\}.
 \label{eq:objetivo-wsp}
\end{equation}
\noindent Todos os vértices têm o mesmo custo; quando as células possuem
áreas uniformes, a contagem aproxima a área queimada. A representação é
determinística, com tempos fixos durante o horizonte, e não incorpora toda
a incerteza ambiental ou as restrições operacionais de incêndios reais.

Embora os tempos de chegada possam ser calculados por caminhos mínimos
para uma alocação fixa, escolher as proteções é combinatoriamente difícil.
Cada decisão altera a propagação e a viabilidade das decisões seguintes.
\citeonline{DelazeriRitt2026} demonstram que a versão de decisão do WSP
é NP-completa, mesmo no caso em que todos os recursos são liberados
simultaneamente. Assim, a dificuldade do problema não se deve apenas à
variação das liberações, e não se espera um algoritmo exato eficiente para
todas as instâncias. Entretanto, a complexidade de pior caso não implica
que toda instância seja igualmente difícil: a estrutura do grafo, os
recursos, os atrasos e as liberações afetam o esforço de solução. Essa
variabilidade motiva a análise das características que tornam uma instância
mais ou menos custosa de resolver.

Essa heterogeneidade justifica diversificar os testes. Os benchmarks
clássicos tornaram-se pequenos e pouco discriminantes para métodos atuais.
\citeonline{DelazeriRitt2026} propuseram um gerador fundamentado no modelo
de propagação superficial de \citeonline{Rothermel1972}, com controle de
tamanho, relevo, vento e configuração dos recursos. Os casos produzidos
permitem estudar essas combinações em condições computacionais controladas.

O WSP se relaciona ao \textit{firefighter problem}, proposto por Hartnell
em 1995, conforme \citeonline{FinbowMacGillivray2009}. Nesse problema, o
incêndio começa em um ou mais vértices de um grafo e evolui em passos
discretos: a cada passo, até $d$ vértices ainda não afetados são defendidos
e, em seguida, o fogo se propaga de cada vértice queimado a todos os seus
vizinhos não defendidos. Vértices queimados ou defendidos mantêm esse estado,
e o processo termina quando o fogo não pode mais avançar. Entre os objetivos
usuais estão maximizar o número de vértices salvos, o que equivale a
minimizar os queimados, e minimizar o tempo até o fim do processo
\cite{FinbowMacGillivray2009}. O WSP, por sua vez, é definido sobre um grafo
direcionado cujos arcos têm tempos de propagação reais não negativos, que
determinam o instante de chegada do fogo a cada vértice a partir da ignição
\cite{DelazeriRitt2026}. Cada recurso de supressão tem um instante de
liberação e só pode ser alocado a um vértice que ainda não queimou nesse
instante, com no máximo um recurso por vértice. O recurso, porém, não bloqueia
o vértice: acrescenta um atraso finito aos arcos de saída, e o próprio
vértice protegido ainda pode queimar. O objetivo é minimizar o número de
vértices queimados antes de um horizonte $H$. Assim, o atraso finito, os
tempos de propagação, as liberações e o horizonte do WSP impedem transferir
diretamente soluções e resultados de complexidade entre as formulações
\cite{DelazeriRitt2026,FinbowMacGillivray2009}.

\subsection{Formulação em programação inteira mista}\label{sec:modelo-mip}

A formulação de \citeonline{DelazeriRitt2026} agrupa os recursos pelos
instantes distintos de liberação $\tau_1<\cdots<\tau_T$. Seja $q_i$ a
quantidade liberada em $\tau_i$, de modo que $\sum_{i=1}^{T}q_i=k$.
A variável binária $x_{iv}$ indica a alocação de um recurso de $\tau_i$
ao vértice $v$; $y_v$ indica que $v$ queima antes de $H>0$; e $a_v$ é
uma variável contínua não negativa associada à chegada do fogo. Com o
mesmo atraso $\Delta$ para os recursos da instância, o modelo é
\begin{align}
 \min\quad &\sum_{v\in V}y_v \label{eq:mip-objetivo}\\
 \text{sujeito a}\quad &a_s=0, \label{eq:mip-ignicao}\\
 &a_v\leq a_u+t_{uv}+\Delta\sum_{i=1}^{T}x_{iu},
       &&(u,v)\in A, \label{eq:mip-propagacao}\\
 &\sum_{v\in V}x_{iv}\leq q_i,
       &&i=1,\ldots,T, \label{eq:mip-capacidade}\\
 &\sum_{i=1}^{T}x_{iv}\leq1,
       &&v\in V, \label{eq:mip-unicidade}\\
 &a_v\geq\tau_i x_{iv},
       &&i=1,\ldots,T,\ v\in V, \label{eq:mip-temporal}\\
 &y_v\geq1-a_v/H,
       &&v\in V, \label{eq:mip-queima}\\
 &x_{iv},y_v\in\{0,1\},\quad a_v\geq0.
       \label{eq:mip-dominios}
\end{align}

A restrição~\eqref{eq:mip-ignicao} fixa a ignição, e
\eqref{eq:mip-propagacao} limita a chegada a cada sucessor pelo tempo
do predecessor mais o deslocamento e o eventual atraso. As restrições
\eqref{eq:mip-capacidade} e \eqref{eq:mip-unicidade} controlam a
capacidade por liberação e impedem duas alocações no mesmo vértice.
A condição~\eqref{eq:mip-temporal} exige que o recurso esteja disponível
antes ou no instante da chegada, quando $x_{iv}=1$.
Por fim, \eqref{eq:mip-queima} obriga $y_v=1$ se $a_v<H$.

Na definição por caminhos mínimos, a chegada é uma igualdade; no modelo,
as desigualdades de propagação apenas a limitam superiormente. Para uma
alocação fixa, os tempos reais de caminhos mínimos satisfazem essas
restrições, e qualquer solução do modelo tem $a_v$ no máximo igual a
esses tempos nos vértices alcançáveis. O objetivo favorece valores que
permitam reduzir $y_v$, preservando a contagem ótima de queima.
Isso não exige que todos os valores de $a_v$ retornados pelo
\textit{solver} sejam tempos exatos, especialmente após o horizonte.
Por isso, as soluções são reavaliadas por caminhos mínimos neste estudo.

Para $n=|V|$ e $m=|A|$, há $nT+n$ variáveis binárias e $n$ contínuas,
totalizando $n(T+2)$. Excluídos os domínios, há
$m+nT+2n+T+1$ restrições: a ordem é $O(nT+n)$ para variáveis e
$O(m+nT+n+T)$ para restrições. Trata-se de uma formulação compacta,
o que não implica tempo polinomial para sua resolução inteira.

\section{Métodos de solução para o WSP}\label{sec:metodos-wsp}

A comparação reúne métodos exatos (MIP e LBBD), heurísticos (ILS e IBS),
uma matheurística (CUTH) e uma linha de base aleatória (RS).
O Quadro~\ref{quad:metodos-wsp} apresenta suas características.
As subseções explicam os mecanismos e os parâmetros; os valores efetivamente
utilizados estão no Quadro~\ref{qua:met-algoritmos}.

\begin{quadro}[htbp]
\centering
\caption{Abordagens consideradas para o WSP.}\label{quad:metodos-wsp}
\begin{singlespace}
\begin{tabular}{|p{1.4cm}|p{2.4cm}|p{8.4cm}|}
\hline
Método & Natureza & Mecanismo principal\\ \hline
MIP & Exato & Otimização conjunta de alocações e propagação por um \textit{solver}.\\ \hline
LBBD & Exato & Alocação no problema mestre e avaliação da propagação no subproblema, com cortes lógicos.\\ \hline
ILS & Heurístico & Busca local e perturbação de alocações.\\ \hline
IBS & Heurístico & Feixe de alocações parciais ordenadas por avaliação heurística.\\ \hline
CUTH & Matheurístico & Corte estático seguido de alocação gulosa temporal.\\ \hline
RS & Heurístico & Construções aleatórias de alocações viáveis, como linha de base.\\ \hline
\end{tabular}
\par\smallskip
{\fonte{elaboração própria, com base em \citeonline{Alvelos2018}, \citeonline{Delazeri2025}, \citeonline{DelazeriRitt2024}, \citeonline{HarrisForbesTaimre2023} e \citeonline{MendesAlvelos2023}.}}
\end{singlespace}
\end{quadro}

\subsection{Programação Inteira Mista}\label{sec:metodo-mip}

A MIP combina decisões discretas de alocação e queima com variáveis
contínuas de chegada, permitindo utilizar um \textit{solver} gratuito ou
comercial de otimização inteira. O modelo inicial de
\citeonline{Alvelos2018} emprega dualidade de programação linear para
impor tempos exatos de caminhos mínimos. A formulação de
\citeonline{DelazeriRitt2026}, apresentada na Seção~\ref{sec:modelo-mip},
utiliza desigualdades de propagação mais simples, suficientes para a
contagem de vértices queimados no objetivo. Essa diferença reduz o esforço
de modelagem da propagação, sem remover a dificuldade combinatória.

Neste trabalho, MIP designa a resolução dessa segunda formulação pelo
Gurobi. O termo modelo básico indica que uma única constante $\Delta$
vale para todas as alocações de uma instância. Não se trata de um modelo
com defesa permanente nem com atraso dependente do recurso.
As decisões $x_{iv}$ determinam $\Lambda$, e os tempos e a queima resultantes
são conferidos sob o modelo de caminhos mínimos.

Durante a resolução, uma alocação viável fornece um limite superior para
$|B_H^{\Lambda}|$, e a exploração do modelo produz limites inferiores.
A diferença entre a melhor incumbente e um limite inferior global válido
origina o \textit{gap} interno do \textit{solver}. Sua redução a zero,
consideradas as tolerâncias numéricas, permite certificar otimalidade.
Já o \textit{gap} empírico da Seção~\ref{sec:met-analise} compara o método
com a melhor solução observada entre as execuções e não constitui essa prova.

A classificação como exato descreve a garantia do método quando a
resolução é concluída, não assegura uma solução ótima dentro do orçamento.
Se o limite de tempo for atingido, podem restar uma solução viável e um
intervalo de otimalidade ainda aberto. O orçamento e a configuração usados
são apresentados no Quadro~\ref{qua:met-algoritmos}.

\subsection{Decomposição de Benders Baseada em Lógica}\label{sec:lbbd}

A LBBD de \citeonline{HarrisForbesTaimre2023} separa as decisões de alocação
da avaliação da propagação. O problema mestre decide onde e em qual
liberação empregar os recursos, respeitando capacidades e unicidade.
Também mantém variáveis de queima que inicialmente aproximam o efeito
das alocações. O subproblema recebe o conjunto protegido, recalcula os
tempos $a_v^{\Lambda}$ por caminhos mínimos e confronta a propagação real
com a solução proposta pelo mestre.

Quando uma proteção ocorre após a chegada do fogo, o subproblema identifica
a incompatibilidade temporal e acrescenta um corte de viabilidade.
Quando a estimativa de queima é otimista, acrescenta cortes de otimalidade
que obrigam o mestre a reconhecer a perda ou modificar as proteções.
Os cortes são derivados da estrutura dos caminhos responsáveis pela
propagação. O procedimento alterna exploração do mestre e separação de
cortes até comprovar a solução ou atingir o critério de parada.

A implementação empregada inclui uma inicialização heurística. Nessa
fase, as alocações são construídas progressivamente pelas liberações,
utilizando uma versão restrita do mestre e verificações de propagação.
A solução resultante inicializa a etapa exata, que permite reconsiderar
as decisões. Essa inicialização pode consumir parte substancial do tempo;
por isso, as duas fases compartilham o orçamento no experimento, conforme
o Quadro~\ref{qua:met-algoritmos} e a Seção~\ref{sec:met-benchmark}.
Um limite da versão heurística restrita não deve ser interpretado como
limite global do problema completo.

A força dos cortes depende do quanto as proteções alteram os caminhos.
Quando o atraso de um único recurso não basta para deslocar a chegada além
de $H$, várias proteções podem ser necessárias para salvar um vértice,
dificultando a identificação de condições fortes de queima e não queima.
Assim, a decomposição pode exigir muitas verificações e cortes.
A LBBD é exata quando sua etapa completa converge com limites válidos;
como na MIP, essa classificação não garante otimalidade dentro do tempo
permitido. Solução, limite, estado de término e fase executada precisam
ser interpretados conjuntamente.

\subsection{Busca Local Iterada}\label{sec:ils}

A ILS de \citeonline{MendesAlvelos2023} alterna melhoria de alocações
completas e perturbações, buscando linhas de contenção de boa qualidade.
Na implementação utilizada \cite{HarrisForbesTaimre2023}, a construção
inicial percorre as liberações e sorteia o próximo vértice entre os
candidatos viáveis de chegada mais próxima. Repete-se essa construção em
múltiplos reinícios e conserva-se a melhor solução inicial.

A busca local tenta deslocar um recurso para a vizinhança dos demais
vértices protegidos, incluindo vizinhos espaciais diagonais existentes
na grade. Os candidatos são ordenados pela chegada do fogo e limitados
pelo parâmetro de vizinhos. Cada movimento é reavaliado por caminhos
mínimos e só pode ser aceito se preservar todas as alocações temporais.
Passagens de melhoria são repetidas até que não se encontre redução
estrita de $|B_H^{\Lambda}|$ na vizinhança examinada.

A perturbação remove uma alocação com probabilidade $p_1$, adiciona um
recurso disponível com probabilidade $p_2$, quando aplicável, ou reposiciona
recursos em uma vizinhança ampliada. O número máximo de modificações limita
os deslocamentos aceitos; o número máximo de falhas interrompe tentativas
inviáveis consecutivas. Após a perturbação, aplica-se novamente a busca
local. A melhor solução é atualizada quando há melhoria, e um contador
controla as iterações consecutivas sem progresso. O experimento também
impõe o limite global de tempo. Os valores desses parâmetros e dos
reinícios e candidatos iniciais estão no Quadro~\ref{qua:met-algoritmos}.

\begin{quadro}[htbp]
\centering
\caption{Pseudocódigo da ILS: construção, melhoria e perturbação.}
\label{qua:pseudo-ils}
\begin{minipage}{0.95\textwidth}
\begin{singlespace}
\textbf{Entrada:} instância, limites de construção e busca, $p_1$, $p_2$.\\
\textbf{Saída:} melhor alocação viável $\Lambda^\star$ registrada.
\begin{enumerate}
 \item Construa alocações aleatórias viáveis em múltiplos reinícios e retenha a melhor.
 \item Aplique busca local e inicialize $\Lambda^\star$; faça $\ell\leftarrow0$.
 \item Enquanto houver tempo e $\ell$ for menor que o limite sem melhoria:
 \begin{enumerate}
  \item Perturbe $\Lambda^\star$ por remoção, adição ou reposicionamento.
  \item Aplique busca local à alocação perturbada, respeitando a viabilidade.
  \item Se o objetivo melhorar, atualize $\Lambda^\star$ e faça $\ell\leftarrow0$; caso contrário, faça $\ell\leftarrow\ell+1$.
 \end{enumerate}
 \item Retorne $\Lambda^\star$.
\end{enumerate}
{\fonte{elaboração própria a partir de \citeonline{MendesAlvelos2023} e da implementação de \citeonline{HarrisForbesTaimre2023}.}}
\end{singlespace}
\end{minipage}
\end{quadro}

\subsection{Busca em Feixe Iterada}\label{sec:ibs}

O IBS de \citeonline{DelazeriRitt2024} explora uma árvore cuja raiz é a
alocação vazia, cujos níveis acompanham os recursos liberados e cujos nós
representam alocações parciais. A largura $\beta$ limita o número de
alocações mantidas por nível, enquanto $\eta$ limita as expansões
preservadas de cada alocação antes da seleção do feixe seguinte.

A avaliação $h_1(\Lambda)=|B_H^{\Lambda}|$ coincide com o objetivo do WSP.
Nos primeiros níveis, porém, decisões que atrasam o fogo podem manter a
mesma contagem de queimados. Para discriminar essas alocações, utiliza-se
\begin{equation}
 h_2(\Lambda)=\sum_{v\in V}\max\{H-a_v^{\Lambda},0\}.
 \label{eq:ibs-h2}
\end{equation}
\noindent Valores menores indicam chegada agregada mais tardia aos
vértices ainda atingidos antes de $H$; vértices inalcançáveis contribuem
com zero. A ordenação inicia com $h_2$ e passa a $h_1$ no instante
$\hat p$ vezes o tempo de queima livre. O parâmetro $\hat p$ controla,
portanto, quando se passa a priorizar diretamente a contagem final.

Em cada expansão, os candidatos são restringidos a um perímetro definido
por uma janela de chegada posterior à liberação. O parâmetro $z$ amplia
essa janela para incluir vértices que serão alcançados mais tarde.
O fator $c$ multiplica o tamanho do perímetro para determinar as tentativas
de construção, e $p$ controla a preferência por candidatos vizinhos de
vértices já protegidos. Quando uma busca em feixe não melhora a incumbente,
atualiza-se $z\leftarrow(z+1)\bmod z_{\max}$; no experimento,
$z_{\max}=3$ produz o ciclo $0,1,2$. Se houver melhoria, mantém-se $z$.
Os demais valores estão no Quadro~\ref{qua:met-algoritmos}.

\begin{quadro}[htbp]
\centering
\caption{Pseudocódigo do IBS com buscas em feixe sucessivas.}
\label{qua:pseudo-ibs}
\begin{minipage}{0.95\textwidth}
\begin{singlespace}
\textbf{Entrada:} instância, $\beta$, $\eta$, $c$, $p$, $\hat p$, $z_{\max}$ e parada.\\
\textbf{Saída:} melhor alocação viável $\Lambda^\star$.
\begin{enumerate}
 \item Inicialize $\Lambda^\star$ com a alocação vazia e faça $z\leftarrow0$.
 \item Enquanto o critério de parada não for satisfeito:
 \begin{enumerate}
  \item Inicie o feixe com a alocação vazia.
  \item Para cada recurso, em ordem de liberação, expanda as alocações do feixe usando o perímetro $z$, o esforço $c$ e a preferência $p$.
  \item Em cada nível, preserve até $\eta$ expansões por alocação e até $\beta$ alocações no novo feixe, ordenadas por $h_2$ ou $h_1$ conforme $\hat p$.
  \item Ao concluir a busca, registre sua solução representativa $\Lambda$.
  \item Se $\Lambda$ melhorar o objetivo, atualize $\Lambda^\star$; caso contrário, faça $z\leftarrow(z+1)\bmod z_{\max}$.
 \end{enumerate}
 \item Retorne $\Lambda^\star$.
\end{enumerate}
{\fonte{elaboração própria a partir de \citeonline{DelazeriRitt2024,Delazeri2025}.}}
\end{singlespace}
\end{minipage}
\end{quadro}

O IBS é estocástico e não prova otimalidade. A solução devolvida por uma
busca em feixe pode ser pior que a incumbente. A distinção entre a sequência
dessas soluções e a sequência de melhores valores sustenta a análise
por STNs, pois apenas as melhorias não conservam todo o percurso registrado.

\subsection{Heurística Baseada em Cortes}\label{sec:cuth}

A CUTH de \citeonline{Delazeri2025} combina um corte estático de vértices
com uma construção gulosa temporal. Para um instante $\bar t$, calcula-se
$B_{\bar t}^{\Lambda_0}$, o conjunto queimado sem recursos antes desse
instante. Busca-se um corte com até $k$ vértices que minimize quantos
vértices permanecem alcançáveis a partir desse conjunto. A variável
binária $\chi_v$ indica participação no corte e $\omega_v$ indica
alcançabilidade. O submodelo é
\begin{align}
 \min\quad &\sum_{v\in V}\omega_v \label{eq:cuth-objetivo}\\
 \text{sujeito a}\quad &\omega_v=1,\quad\chi_v=0,
       &&v\in B_{\bar t}^{\Lambda_0}, \label{eq:cuth-fontes}\\
 &\omega_u\leq\omega_v+\chi_u,
       &&(u,v)\in A, \label{eq:cuth-propagacao}\\
 &\sum_{v\in V}\chi_v\leq k, \label{eq:cuth-capacidade}\\
 &\chi_v,\omega_v\in\{0,1\},&&v\in V.
\end{align}

Os vértices já queimados são alcançáveis e não podem integrar o corte.
A condição~\eqref{eq:cuth-propagacao} estende a alcançabilidade pelos arcos,
exceto quando o vértice de origem pertence ao corte. O próprio vértice do
corte pode ser alcançado, mas seus sucessores podem deixar de sê-lo.
A Figura~\ref{fig:corte-estatico} ilustra essa separação. O limite de
$k$ proteções usa a capacidade total e ainda não assegura que todos os
recursos estejam disponíveis antes da chegada do fogo.

\begin{figure}[htbp]
\centering
\caption{Exemplo de corte estático utilizado como candidato à contenção.}
\label{fig:corte-estatico}
\begin{singlespace}
\begin{tikzpicture}[v/.style={circle,draw,minimum size=8mm},
 seta/.style={-{Stealth[length=2mm]},thick}]
\node[v,fill=black!15] (s) at (0,0) {$s$};
\node[v,fill=black!15] (b) at (1.8,0) {$b$};
\node[v,fill=blue!15] (c1) at (3.6,0.9) {$c_1$};
\node[v,fill=blue!15] (c2) at (3.6,-0.9) {$c_2$};
\node[v] (d1) at (5.8,0.9) {$d_1$};
\node[v] (d2) at (5.8,-0.9) {$d_2$};
\draw[seta] (s)--(b);
\draw[seta] (b)--(c1);\draw[seta] (b)--(c2);
\draw[seta,dashed] (c1)--(d1);\draw[seta,dashed] (c2)--(d2);
\node[font=\small] at (0.9,-1.8) {já queimados};
\node[font=\small] at (3.6,-1.8) {corte};
\node[font=\small] at (5.8,-1.8) {separados};
\end{tikzpicture}
\par\smallskip
{\small Nota: cinza indica os vértices já queimados em $\bar t$; azul, o corte
$P=\{c_1,c_2\}$; branco, a região separada. Os arcos tracejados deixam de
propagar alcançabilidade no submodelo. No WSP, a proteção produz atraso
finito e depende da alocação temporal dos recursos.\par}
{\fonte{elaboração própria, com base em \citeonline{Delazeri2025}.}}
\end{singlespace}
\end{figure}

Após obter $P=\{v:\chi_v=1\}$, recalculam-se as chegadas considerando
os atrasos em $P$. A construção gulosa associa o recurso de liberação
mais próxima ao vértice de $P$ com menor chegada do fogo, desde que a
liberação não seja posterior à chegada. Repete-se até alocar todo o corte
ou identificar inviabilidade. Uma busca binária sobre $\bar t$ procura um
corte que possa ser implantado. A tolerância temporal, o tempo por chamada
do submodelo e o orçamento global estão no Quadro~\ref{qua:met-algoritmos}.

A CUTH é matheurística porque resolve um submodelo exato dentro de uma
estratégia heurística para o WSP. Os tempos definem o conjunto queimado
e a viabilidade temporal, mas não entram na propagação de alcançabilidade
do corte estático. Seu ótimo não é um limite inferior certificado do WSP;
a alocação viável obtida fornece um limite superior para $|B_H^{\Lambda}|$.

\subsection{Busca Aleatória}\label{sec:rs}

A RS serve de linha de base para avaliar o ganho das estratégias orientadas
\cite{DelazeriRitt2026}. A cada construção, começa com a alocação vazia,
percorre as liberações em ordem e sorteia vértices elegíveis, respeitando
a capacidade e a unicidade. As chegadas são atualizadas à medida que as
proteções alteram a propagação. Ao fim, compara-se o objetivo com a melhor
solução anterior. A construção é repetida até a parada.

\begin{quadro}[htbp]
\centering
\caption{Pseudocódigo da RS como linha de base de desempenho.}
\label{qua:pseudo-rs}
\begin{minipage}{0.95\textwidth}
\begin{singlespace}
\begin{enumerate}
 \item Inicialize $\Lambda^\star$ com a alocação vazia.
 \item Enquanto houver tempo, inicie uma nova alocação vazia $\Lambda$.
 \item Para cada liberação e recurso disponível, recalcule os vértices sem proteção que ainda não queimaram, sorteie um deles e aloque o recurso; se não houver candidato, deixe-o sem uso.
 \item Ao concluir a construção, atualize $\Lambda^\star$ se $|B_H^{\Lambda}|$ for menor; retorne à etapa 2 enquanto houver tempo.
 \item Retorne $\Lambda^\star$.
\end{enumerate}
{\fonte{elaboração própria a partir de \citeonline{DelazeriRitt2026}.}}
\end{singlespace}
\end{minipage}
\end{quadro}

O IBS é o método de referência para o estudo das trajetórias; a RS é o
método de base usado como ponto de comparação de desempenho. A semente
controla a amostragem, e suas soluções viáveis fornecem limites superiores,
sem certificar otimalidade. A configuração da RS está no
Quadro~\ref{qua:met-algoritmos}.

\section{Análise de espaço de instâncias}\label{sec:isa}

Médias de desempenho podem ocultar regiões favoráveis a diferentes
algoritmos. A ISA torna essa heterogeneidade explícita ao
relacionar atributos mensuráveis dos problemas ao desempenho
\cite{MunozEtAl2018,SmithMilesMunoz2023}. Essa abordagem se relaciona
à seleção de algoritmos de \citeonline{Rice1976} e à avaliação objetiva
em diferentes regiões desenvolvida por \citeonline{SmithMilesEtAl2014}.

Cada instância $I_i$ é representada por um vetor
$\phi(I_i)\in\mathbb{R}^m$ de atributos e por medidas de desempenho
$y_{i\alpha}$ dos algoritmos $\alpha$. No WSP, os atributos descrevem
tamanho, tempos, recursos, queima livre e topologia. O desempenho pode
expressar qualidade ou tempo; quando há repetições estocásticas, deve ser
agregado por instância, preservando também sua dispersão. A definição do
\textit{gap} adotado neste trabalho está na Seção~\ref{sec:met-analise}.

Uma projeção $\psi:\mathbb{R}^m\rightarrow\mathbb{R}^2$ associa cada
instância a $z_i=\psi(\phi(I_i))$. Proximidade no mapa expressa similaridade
nos atributos após as transformações utilizadas. A sobreposição do desempenho permite
examinar agrupamentos, lacunas e \textit{footprints}: regiões associadas
ao critério de bom desempenho. Na configuração deste estudo, essas
regiões são inferidas a partir de predições, não medidas em todos os pontos.
Sua interpretação depende da coleção, dos atributos e dos modelos.

\subsection{O fluxo MATILDA}\label{sec:matilda}

A \textit{Melbourne Algorithm Test Instance Library with Data Analytics}
(MATILDA) organiza a ISA em preparação, seleção, projeção,
predição e análise de regiões \cite{SmithMilesMunoz2023}. As etapas
possuem os seguintes papéis:
\begin{alineas}
 \item \textbf{PRELIM}: verifica e transforma os dados para tratar assimetrias e escalas;
 \item \textbf{SIFTED}: seleciona atributos informativos e reduz redundâncias;
 \item \textbf{PILOT}: projeta atributos e tendências de desempenho em duas dimensões;
 \item \textbf{CLOISTER}: estima limites plausíveis a partir das relações entre atributos;
 \item \textbf{PYTHIA}: ajusta classificadores do critério de bom desempenho;
 \item \textbf{TRACE}: constrói e descreve os footprints segundo as classes utilizadas.
\end{alineas}

A análise de componentes principais (\textit{Principal Component Analysis},
PCA) teve finalidade exploratória, enquanto PILOT utiliza informação de
desempenho. No pyInstanceSpace, PYTHIA e CLOISTER dependem de PILOT,
mas não uma da outra. TRACE utiliza as predições de PYTHIA, sem depender
diretamente da saída de CLOISTER \cite{PyInstanceSpace021}.
A Seção~\ref{sec:met-matilda} apresenta a configuração adotada.

Transformações e seleção afetam a geometria; o limiar de bom desempenho
afeta os rótulos; e poucas observações ou classes desbalanceadas podem
tornar os classificadores pouco informativos. Por isso, mapas e regiões
precisam ser acompanhados de diagnósticos. Avaliação na própria coleção
não demonstra generalização para novas instâncias.

\subsection{Geração orientada pelo espaço de instâncias}

Lacunas no mapa podem refletir combinações inviáveis de atributos ou
insuficiência dos testes existentes. \citeonline{SmithMilesBowly2015}
propõem investigar a segunda possibilidade definindo pontos-alvo e
utilizando busca evolutiva para gerar instâncias próximas a eles.
Para uma candidata $I(\theta)$, produzida pelos parâmetros $\theta$,
um objetivo de geração é
\begin{equation}
 \min_{\theta}\left\|\psi\bigl(\phi(I(\theta))\bigr)-z^{\star}\right\|_2,
 \label{eq:geracao-alvo-isa}
\end{equation}
\noindent sendo $z^{\star}$ o alvo. Seleção, recombinação e mutação
exploram configurações do gerador sujeitas às restrições do domínio.

A projeção das candidatas deve reutilizar as transformações e atributos
do mapa que definiu os alvos. Mantida essa referência e verificada a validade
das instâncias, a redução das distâncias aos pontos do domínio demonstra
maior cobertura da representação. A dificuldade dos casos e sua contribuição
à discriminação de algoritmos são avaliadas pelo desempenho medido.
Um reajuste da ISA define outra geometria. A cobertura é medida no mapa
que orientou a geração; o mapa reajustado permite analisar as relações
entre atributos e desempenho na coleção ampliada.

\section{Caracterização topológica de grafos direcionados}\label{sec:topologia}

O tamanho e as estatísticas dos pesos não determinam a organização dos caminhos
de um grafo. Medidas de redes complementam esses descritores ao representar
conectividade local, influência e acessibilidade \cite{Newman2018}.
A semântica dos pesos é essencial: $t_{uv}$ é uma distância temporal,
adequada às medidas de caminhos mínimos. Para medidas em que peso significa
intensidade da conexão, emprega-se $w_{uv}=1/t_{uv}$, fazendo arcos rápidos
corresponderem a ligações fortes.

Os graus de entrada e saída contam, respectivamente, os arcos que chegam
a um vértice e os que dele partem. Razões e diferenças normalizadas
descrevem a assimetria entre recepção e disseminação. Em grades regulares,
fronteiras e remoção de regiões inalcançáveis introduzem variação nessas
medidas. O PageRank resume influência global de entrada por um passeio
aleatório com teletransporte e tratamento de vértices sem saída
\cite{BrinPage1998}. A Busca de Tópicos Induzida por Hiperligações (HITS, do inglês:
\textit{hyperlink-induced topic search}) distingue \textit{authorities}, que recebem
ligações de bons \textit{hubs}, e \textit{hubs}, que apontam para boas \textit{authorities};
seus escores derivam dos vetores singulares dominantes da matriz de
adjacência ou de forças \cite{Kleinberg1999}.

A intermediação quantifica a participação de um vértice nos caminhos
mínimos entre pares:
\begin{equation}
 C_B(v)=\sum_{\substack{p\neq q\\p,q\neq v}}
 \frac{\sigma_{pq}(v)}{\sigma_{pq}},
 \label{eq:betweenness}
\end{equation}
\noindent em que $\sigma_{pq}$ conta os caminhos mínimos de $p$ a $q$,
e $\sigma_{pq}(v)$, aqueles que atravessam $v$. Pares desconectados
contribuem com zero. Em um dígrafo com $n>2$, a normalização por
$(n-1)(n-2)$ considera pares ordenados. Alta intermediação pode sinalizar
gargalos, mas a medida agrega pares diversos, enquanto o fogo parte de
uma ignição específica.

A centralidade harmônica mede acessibilidade mesmo sem conexão forte:
\begin{equation}
 C_H^{\mathrm{out}}(v)=\sum_{u\neq v}\frac{1}{d_G(v,u)},\qquad
 C_H^{\mathrm{in}}(v)=\sum_{u\neq v}\frac{1}{d_G(u,v)},
 \label{eq:centralidade-harmonica}
\end{equation}
\noindent Vértices inalcançáveis têm distância infinita e contribuem
com zero para a soma. As formas de saída e entrada expressam
facilidade de alcançar a rede e de ser alcançado. Com tempos como distância,
incorporam a rapidez da propagação; a normalização aplicada às instâncias
está na Seção~\ref{sec:met-topologia}.

O agrupamento direcionado mede fechamento de vizinhanças por triângulos.
A variante binária de \citeonline{Fagiolo2007} considera suas diferentes
orientações. Em grades sem diagonais, a ausência de triângulos pode tornar
a medida constante. Essa propriedade ilustra por que uma métrica de rede
não deve ser considerada informativa antes de examinar os grafos utilizados.

As medidas são nodais, mas a ISA utiliza uma observação por instância.
Por isso, suas distribuições são resumidas por média, dispersão, quantis e
extremos. A agregação conserva aspectos do perfil estrutural, mas perde
localização dos vértices e pode introduzir redundância com tamanho e pesos.
Filtros de variância, correlação e estabilidade numérica são necessários.

O benefício de proteger um vértice depende da ignição, do horizonte, das
liberações e da interação entre alocações. As centralidades fornecem
hipóteses sobre sua importância estrutural, que podem ser examinadas
na representação das instâncias e nos diagnósticos de desempenho.

\section{Redes de trajetórias de busca}\label{sec:stn}

As STNs representam percursos de busca por grafos direcionados
\cite{OchoaMalanBlum2021}. Enquanto uma curva de convergência mostra quando
a melhor qualidade muda, a rede permite examinar se execuções compartilham
estados e transições, retornam a soluções anteriores ou permanecem dispersas.

Uma solução representativa resume a busca em um instante de registro,
como a incumbente de um método de ponto único ou a melhor solução de uma
população. Uma localização é uma parte não vazia de uma partição do espaço
de soluções: pode corresponder a uma solução ou agrupar soluções consideradas
equivalentes. A trajetória é a sequência de localizações registradas.
A STN $\mathcal{G}_{\mathrm{STN}}=(N,E)$ reúne essas localizações como nós
e as transições consecutivas como arestas, ponderadas por sua frequência.
Atributos dos nós podem indicar qualidade, frequência de visita e condição
inicial ou terminal. Trajetórias de várias execuções podem ser agregadas
quando compartilham a mesma definição de localização
\cite{ChaconSartoriEtAl2023,OchoaMalanBlum2021}.

A codificação e a frequência de registro determinam o que a rede representa.
Uma partição muito fina reduz sobreposições; uma partição grosseira pode
fundir estados relevantes para a análise. No WSP, o instante também faz
parte da decisão de proteção de um vértice. Registrar apenas
melhorias também produz uma rede diferente daquela formada por todas as
soluções devolvidas pelas iterações. O protocolo específico deste trabalho
é definido na Seção~\ref{sec:met-trajetorias}.

Entre os descritores estão o número de nós e arestas, os nós de melhor
avaliação observada ($n_{\mathrm{best}}$), os terminais fora desse conjunto
($n_{\mathrm{end}}$) e o compartilhamento entre execuções. A
\textit{best-strength} soma os pesos das arestas que chegam aos melhores
nós e divide pelo número de execuções. Transições repetidas podem fazer
essa razão exceder um: trata-se de frequência normalizada de transições.

Nenhuma dessas métricas determina isoladamente a qualidade da busca.
Uma rede pequena pode representar convergência ou estagnação; uma rede
grande, diversificação produtiva ou dispersão. Qualidade, recorrência,
compartilhamento e evolução precisam ser interpretados conjuntamente.
Assim, a topologia descreve o grafo do problema, a ISA situa a instância
em relação às demais e a STN descreve as soluções visitadas por um algoritmo. Neste trabalho,
ela é aplicada ao IBS.
